%! AP Statistics — Unit 4, Lesson 4.3 — Guided Notes KEY
\documentclass[10pt]{article}
\usepackage{apstats-article}
\usepackage{apstats-key}

\begin{document}

\pageheader{Unit 4, Lesson 4.3}{Guided Notes --- KEY}
\namedateperiod

\vspace{0.12in}

\begin{objectivebox}
\small By the end of this lesson, I will be able to\dots\\[4pt]
\ans{Calculate the probability of an event using sample spaces and the classical definition.}\\[1pt]
\ans{Apply the complement rule to find $P(E^c) = 1 - P(E)$.}\\[1pt]
\ans{Interpret a probability as the long-run relative frequency of an event.}
\end{objectivebox}

\vspace{0.1in}

\begin{vocabbox}
\termblanklong{Sample space ($S$)} \ans{The set of all possible non-overlapping outcomes of a random process.}
\termblanklong{Event} \ans{A subset of the sample space --- one or more outcomes.}
\termblanklong{Equally likely outcomes} \ans{Each outcome in the sample space has the same probability of occurring.}
\termblanklong{Classical probability} \ans{$P(E) = \dfrac{\text{number of outcomes in }E}{\text{total number of outcomes in }S}$, when all outcomes are equally likely.}
\termblanklong{Complement ($E^c$)} \ans{The event that $E$ does NOT occur; consists of all outcomes in $S$ that are not in $E$.}
\termblanklong{Complement rule} \ans{$P(E^c) = 1 - P(E)$.}
\end{vocabbox}

\vspace{0.1in}

\begin{hookbox}
\small A bag holds 3 red and 5 blue poker chips.
\begin{itemize}[itemsep=3pt]
  \item $S = \{$ \ans{R1, R2, R3, B1, B2, B3, B4, B5 (or simply: red, blue treated as 8 equally likely chips)} $\}$
  \item Total outcomes: \ans{8}
  \item Outcomes that are ``red'': \ans{3}
  \item $P(\text{red}) = \ans{3/8} = \ans{0.375}$
\end{itemize}
\end{hookbox}

\vspace{0.1in}

\begin{notesbox}{The Classical Definition of Probability}
\small
\textbf{Sample Space:} The sample space is the set of all possible \ans{non-overlapping} outcomes
(no outcome counted \ans{more than once}).

\vspace{8pt}
\textbf{Classical Probability Formula:}
\[
  P(E) = \frac{\ans{\text{number of outcomes in }E}}{\ans{\text{total number of equally likely outcomes in }S}}
\]

\vspace{6pt}
\textbf{Probability Range:} $\ans{0} \le P(E) \le \ans{1}$

\vspace{10pt}
\textbf{Worked Example:} $S = \{1, 2, 3, 4, 5, 6\}$
\vspace{4pt}
\renewcommand{\arraystretch}{1.6}
\begin{tabularx}{\linewidth}{@{} l X @{}}
Event & Probability \\
\hline
$E_1$: roll a 5 & \ans{$P(E_1) = 1/6 \approx 0.167$} \\
$E_2$: roll an even number & \ans{$P(E_2) = 3/6 = 1/2 = 0.5$} \\
$E_3$: roll a number greater than 4 & \ans{$P(E_3) = 2/6 = 1/3 \approx 0.333$} \\
\end{tabularx}
\end{notesbox}

\vspace{0.1in}

\begin{notesbox}{The Complement Rule}
\small
The complement of $E$ is the event that $E$ \ans{does not occur}.\\[6pt]
\textbf{Complement Rule:} $P(E^c) = \ans{1 - P(E)}$\\[6pt]
$P(E) + P(E^c) = \ans{1}$ because $E$ and $E^c$ together include \ans{every outcome} in $S$.

\vspace{10pt}
\begin{itemize}[itemsep=4pt]
  \item $P(\text{NOT a 5}) = \ans{1 - 1/6 = 5/6}$
  \item $P(\text{odd}) = \ans{1 - 1/2 = 1/2}$
  \item $P(\text{roll} \le 4) = \ans{1 - 1/3 = 2/3}$
\end{itemize}
\end{notesbox}

\vspace{0.1in}

\begin{notesbox}{Interpreting Probability --- Long Run}
\small
\textbf{VAR-4.B.1:} Probabilities of events in \ans{repeatable} situations can be
interpreted as the relative frequency with which the event will occur \ans{in the long run}.

\vspace{8pt}
``If this random process is repeated many times, the event \blank{4cm} will occur in
approximately \blank{3cm} of the repetitions.''

\vspace{10pt}
\begin{enumerate}[leftmargin=1.5em, itemsep=10pt]
  \item \ans{If a fair die is rolled many times, about 1 out of every 6 rolls will land on 5.}
  \item \ans{If we observe many days with similar weather conditions, about 40\% of them will have rain.}
\end{enumerate}
\end{notesbox}

\vspace{0.1in}

\begin{practicebox}
\small

\vspace{4pt}
\begin{enumerate}[leftmargin=1.5em, itemsep=8pt]
  \item \ans{36 total outcomes (6 $\times$ 6).}
  \item \ans{Sum = 7: (1,6),(2,5),(3,4),(4,3),(5,2),(6,1) → 6 outcomes. $P = 6/36 = 1/6$.}
  \item \ans{Sum $\geq$ 10: (4,6),(5,5),(6,4),(5,6),(6,5),(6,6) → 6 outcomes. $P = 6/36 = 1/6$.}
  \item \ans{$P(\text{NOT sum} \geq 10) = 1 - 1/6 = 5/6$.}
  \item \ans{If two fair dice are rolled many times, about 1 in every 6 rolls will produce a sum of 7.}
\end{enumerate}
\end{practicebox}

\end{document}
